\documentclass[11pt]{article}

\usepackage[margin=1in]{geometry}
\usepackage[T1]{fontenc}
\usepackage[utf8]{inputenc}
\usepackage{lmodern}
\usepackage{enumitem}
\usepackage{hyperref}

\hypersetup{
  colorlinks=true,
  linkcolor=blue!50!black,
  urlcolor=blue!50!black,
  citecolor=blue!50!black,
  pdftitle={Discovery Lens Artifacts for p-vs-np}
}

\title{Discovery Lens Artifacts for \texttt{p-vs-np}}
\author{Neuro-symbolic AI Group\\Idiap Research Institute}
\date{March 17, 2026}

\begin{document}
\maketitle

\begin{abstract}
This report records the structured discovery artifacts produced by the category-theoretic and HoTT reframing lenses. These lenses complement the outer proof-discovery language by exposing composition, mediation, equivalence, transport, and coherence obligations explicitly.
\end{abstract}

\section{Summary}
\begin{itemize}[leftmargin=*,itemsep=0.2em]
\item Category theory is used here as a structural/compositional discovery lens for composition, mediation, commutativity, and universal properties.
\item HoTT is used here as an equivalence/coherence discovery lens for sameness, transport, and representation coherence.
\item Neither lens replaces PDL. Each enriches the outer discovery logic with a specific mode of material inference.
\end{itemize}

\section{categorical-reframing}
\textbf{Prompt title.} Categorical Reframing Prompt

\textbf{Epistemic lens.} category-theoretic

\textbf{Required move.} Recast the blocked frontier as a diagrammatic composition problem: identify objects, morphisms, mediating objects, and the exact commutativity or factorization obligation.

\textbf{Linked queue IDs.} proofq.p-vs-np.001, proofq.p-vs-np.002, proofq.p-vs-np.005, proofq.p-vs-np.009, proofq.p-vs-np.010, proofq.p-vs-np.013, proofq.p-vs-np.018, proofq.p-vs-np.020, proofq.p-vs-np.021, proofq.p-vs-np.006, proofq.p-vs-np.007, proofq.p-vs-np.008

\textbf{Linked breakthrough families.} realization, transport, bridge-specificity, barrier-aware transfer, alignment, constructive instantiation, semantic bridge

\subsection*{Lens Artifacts}
\begin{itemize}[leftmargin=*,itemsep=0.2em]
\item \textbf{Diagram candidates}: Diagram candidate: local route for 'carrier-realization' should factor from source-side object to mediating carrier implementing 'An explicit obstruction family, witness family, surrogate promise problem, or quantitative measure that is realized from the source theorem rather than only named by role.' and then into the target activation interface.; Diagram candidate: local route for 'formula-size-transport' should factor from source-side object to mediating carrier implementing 'A transport law, compilation theorem, or complexity measure that is stable under the source-to-target model shift.' and then into the target activation interface.; Diagram candidate: local route for 'target-grounding' should factor from source-side object to mediating carrier implementing 'No new target theorem; discovery should instead seek the carrier, transfer law, or alignment theorem that makes the target theorem applicable.' and then into the target activation interface.
\item \textbf{Commutativity obligations}: Commutativity obligation: the route for 'carrier-realization' must commute between source theorem hypotheses, the mediating carrier, and the target theorem hypotheses.; Commutativity obligation: the route for 'formula-size-transport' must commute between source theorem hypotheses, the mediating carrier, and the target theorem hypotheses.; Commutativity obligation: the route for 'target-grounding' must commute between source theorem hypotheses, the mediating carrier, and the target theorem hypotheses.
\item \textbf{Universal property candidates}: Universal-property candidate: search for an intermediate object universal for mediating 'Turn a typed but schematic carrier into an explicit mathematical object family with admissibility conditions and a forcing theorem.'.; Universal-property candidate: search for an intermediate object universal for mediating 'Find a theorem that preserves the relevant invariant across a cross-model move into formula-size.'.; Universal-property candidate: search for an intermediate object universal for mediating 'Stop treating the target theorem as the open problem and relocate search to the bridge that would activate it.'.
\item \textbf{Negative outputs}: non\_commuting\_route; missing\_universal\_object; non\_functorial\_transport; broken\_factorization
\end{itemize}
\section{equivalence-coherence-reframing}
\textbf{Prompt title.} Equivalence / Coherence Reframing Prompt

\textbf{Epistemic lens.} hott

\textbf{Required move.} Recast the blocked frontier as an identity, equivalence, transport, or coherence problem rather than as a missing brute-force bridge.

\textbf{Linked queue IDs.} proofq.p-vs-np.001, proofq.p-vs-np.002, proofq.p-vs-np.005, proofq.p-vs-np.009, proofq.p-vs-np.010, proofq.p-vs-np.013, proofq.p-vs-np.018, proofq.p-vs-np.020, proofq.p-vs-np.021, proofq.p-vs-np.003, proofq.p-vs-np.011, proofq.p-vs-np.023

\textbf{Linked breakthrough families.} realization, transport, alignment, explicitness, semantic bridge

\subsection*{Lens Artifacts}
\begin{itemize}[leftmargin=*,itemsep=0.2em]
\item \textbf{Equivalence claims}: Equivalence claim: determine whether the live carrier for 'carrier-realization' is identical to, equivalent to, or only analogous with the target-side object class.; Equivalence claim: determine whether the live carrier for 'formula-size-transport' is identical to, equivalent to, or only analogous with the target-side object class.; Equivalence claim: determine whether the live carrier for 'target-alignment' is identical to, equivalent to, or only analogous with the target-side object class.
\item \textbf{Transport obligations}: Transport obligation: move the invariant across the proposed equivalence for 'carrier-realization' without changing the target claim.; Transport obligation: move the invariant across the proposed equivalence for 'formula-size-transport' without changing the target claim.; Transport obligation: move the invariant across the proposed equivalence for 'target-alignment' without changing the target claim.
\item \textbf{Coherence failures}: Coherence failure probe: check whether different identifications of the carrier for 'carrier-realization' induce the same transported statement.; Coherence failure probe: check whether different identifications of the carrier for 'formula-size-transport' induce the same transported statement.; Coherence failure probe: check whether different identifications of the carrier for 'target-alignment' induce the same transported statement.
\item \textbf{Negative outputs}: no\_equivalence\_witness; transport\_incoherent; weak\_equivalence\_only; identity\_mismatch
\end{itemize}

\end{document}
